% 原书第 122—126 页；6.9。
\section{潜在结局}
\label{sec:ch06-potential-outcomes}

下面介绍另一种不以结构因果模型为基础的因果推断方法。这一框架通常称为\textbf{潜在结局（potential outcomes）}框架或\textbf{Rubin 因果模型（Rubin causal model）}，在社会科学中得到广泛使用。其思想可追溯到 Neyman［1923］和 Fisher［1925］，二者主要讨论随机化实验。Rubin［1974］将这些思想扩展到观测性研究。更详细的介绍见 Rubin［2005］、Morgan and Winship［2007］和 Imbens and Rubin［2015］。

\subsection{定义与例子}

为解释潜在结局，重新考察例 3.4 中的眼科医生，并用这一框架重新表述。现在不从随机变量出发，而从 $n$ 名患者或单位 $u=1,\ldots,n$ 出发，每个单位可能接受，也可能不接受治疗。为每位患者 $u$ 指定两个潜在结局：$B_u(t=1)$ 表示患者接受治疗（$T=1$）时，会失明（$B=1$）还是痊愈（$B=0$）；类似地，$B_u(t=0)$ 表示不治疗（$T=0$）时会发生什么。两个潜在结局都被假定为确定性的。对每位患者，治疗要么有效，要么无效，其中不包含随机性。如果 $B_u(t=1)=0$ 且 $B_u(t=0)=1$，就说治疗对单位 $u$ 具有正面的作用。

但实践中无法核查这些条件。“\textbf{因果推断的基本问题（fundamental problem of causal inference）}”［Holland, 1986］指出：对每个单位 $u$，只能观测 $B_u(t=1)$ 或 $B_u(t=0)$ 中的一个，不能同时观测二者。原因是，一旦决定对某人治疗，就无法回到过去撤销治疗。反方向也一样：如果决定不治疗，之后虽然仍可给予治疗，但那时的结果已不能解释为变量 $B_u(t=1)$ 的一个结局，例如患者可能已在此期间自行康复。因此，我们只能观测其中一个潜在结局；另一个未观测量则成为反事实结局。

表 6.2 给出了上例的一个假想数据集；事实上，这些数据按例 3.4 的模型抽样得到。为使表 6.2 的表示成立，通常隐含假定\textbf{稳定单位治疗值假设（stable unit treatment value assumption，SUTVA）}［Rubin, 2005］。它要求单位之间不存在干扰，例如某个单位的潜在结局不依赖其他单位接受了何种治疗［Cox, 1958］；还要求潜在结局不依赖于该治疗是如何或为何被给予的。第 6.9.2 节将说明，若数据由结构因果模型按本例的方式生成，SUTVA 就成立。

潜在结局告诉我们治疗在个体层面的效应。定义\textbf{单位层面的因果效应（unit-level causal effect）}为 $B_u(t=1)-B_u(t=0)$，平均因果效应为
\begin{equation}
 \mathrm{CE}=\frac1n\sum_{u=1}^n
 \bigl(B_u(t=1)-B_u(t=0)\bigr).
 \tag{6.25}
\end{equation}

\begin{table}[htbp]
\centering
\begin{tikzpicture}[x=1cm,y=1cm,font=\small]
 \draw (0,0) rectangle (11.2,.9);
 \foreach \xx in {.7,1.7,4.3,6.9}{\draw (\xx,0)--(\xx,.9);}
 \node[align=center] at (.35,.45) {单位\\$u$};
 \node[align=center] at (1.2,.45) {治疗\\$T$};
 \node[align=center] at (3,.45) {潜在结局\\$B_u(t=0)$};
 \node[align=center] at (5.6,.45) {潜在结局\\$B_u(t=1)$};
 \node[align=center] at (9.05,.45) {单位层面的因果效应\\$B_u(t=1)-B_u(t=0)$};
 \foreach \u/\tt/\bz/\bo/\ee/\rr in
 {1/1/1/0/-1/1,2/0/1/0/-1/2,3/1/1/0/-1/3,
 43/1/1/0/-1/5,44/0/0/1/1/6,45/0/1/0/-1/7,
 119/1/1/0/-1/9,120/1/0/1/1/10,121/0/1/0/-1/11,
 200/0/1/0/-1/13}{
  \pgfmathsetmacro{\yy}{-.45*\rr}
  \fill[gray!25] (.7,\yy) rectangle (1.7,\yy+.45);
  \ifnum\tt=1
   \fill[gray!25] (4.3,\yy) rectangle (6.9,\yy+.45);
  \else
   \fill[gray!25] (1.7,\yy) rectangle (4.3,\yy+.45);
  \fi
  \draw (0,\yy) rectangle (11.2,\yy+.45);
  \foreach \xx in {.7,1.7,4.3,6.9}{\draw (\xx,\yy)--(\xx,\yy+.45);}
  \node at (.35,\yy+.225) {$\u$};
  \node at (1.2,\yy+.225) {$\tt$};
  \node at (3,\yy+.225) {$\bz$};
  \node at (5.6,\yy+.225) {$\bo$};
  \node at (9.05,\yy+.225) {$\ee$};
 }
 \foreach \rr in {4,8,12}{
  \pgfmathsetmacro{\yy}{-.45*\rr}
  \draw (0,\yy) rectangle (11.2,\yy+.45);
  \draw (.7,\yy)--(.7,\yy+.45);
  \node at (.35,\yy+.225) {$\vdots$};
 }
\end{tikzpicture}
\caption{用潜在结局表示例 3.4。对于每位患者或单位，两个潜在结局中只能观测一个。灰底表示观测到的信息。治疗 $T$ 对几乎所有患者都有益；只有 200 例中的 2 例，治疗会伤害患者并导致失明（$B=1$）。虽然对大多数患者给予治疗（$T=1$）是好选择，但对患者 $u=120$，这恰恰是错误决定。}
\end{table}


因果推断的基本问题，使我们无法直接计算式 (6.25)。假设在一个完全随机化实验中，单位 $u\in U_0\subset\{1,\ldots,n\}$ 接受 $T=0$，单位 $u\in U_1=U_0^{\mathsf C}$ 接受 $T=1$。Neyman［1923］证明，
\begin{equation}
 \widehat{\mathrm{CE}}
 :=\frac1{\#U_0}\sum_{u\in U_0}B_u(t=1)
 -\frac1{\#U_1}\sum_{u\in U_1}B_u(t=0)
 \tag{6.26}
\end{equation}
是式 (6.25) 的无偏估计量。这里，$\widehat{\mathrm{CE}}$ 的随机性来自随机分配，它决定每个单位的两个潜在结局中，哪一个能够被观测；结局本身被视为隐藏的，而非随机的。注意，式 (6.26) 只包含观测量，因此研究完成后即可计算。
\begin{quote}\small
\textbf{译注：}原式 (6.26) 的两组下标与前文定义对调。按前文 $U_0$ 为未治疗组、$U_1$ 为治疗组，原式实际包含两组未观测到的潜在结局，与随后“只包含观测量”的文字矛盾。原印形式在随机分组下仍可无偏，问题在于其不可观测、不能直接计算。与式 (6.25) 一致且可计算的通常差均值估计量应为
\[
 \widehat{\mathrm{CE}}
 =\frac1{\#U_1}\sum_{u\in U_1}B_u(t=1)
 -\frac1{\#U_0}\sum_{u\in U_0}B_u(t=0).
\]
上方编号公式保留原书印法；本译注给出其建议读法。
\end{quote}

关于两种方法中哪一种更适合实际应用，已有广泛争论，例如 Pearl［1995］、Imbens and Rubin［1995］、Rubin［2004］、Lauritzen［2004］。本书不打算积极参与这一争论，而只提及以下三个结果：（1）说明如何将潜在结局表示为反事实［Pearl, 2009，第 3.6.3 节］；（2）两种框架存在逻辑等价性［Galles and Pearl, 1998；Halpern, 2000］；（3）介绍一个近来提出、使两种视角更加接近的框架［Richardson and Robins, 2013］。

\subsection{潜在结局与结构因果模型的关系}

在结构因果模型中，可以用反事实的语言表示潜在结局，见第 6.4 节。在眼科医生的例子中，模型 $\mathfrak C$ 满足 $T=N_T$ 和 $B=T\cdot N_B+(1-T)\cdot(1-N_B)$。因此，可以用 $N_B$ 和 $N_T$ 的特定取值表示每位患者。例如，表 6.2 的患者 43 由 $N_T=1,N_B=0$ 表征，患者 44 则满足 $N_T=0,N_B=1$。$t=0$ 和 $t=1$ 对应于对 $T$ 的干预。归纳起来，
\begin{equation}
 \underbrace{B_u(t=\widetilde t)}_{\text{潜在结局}}
 =\underbrace{\text{模型 }\mathfrak C\mid\mathbf N=\mathbf n_u;\dox{T:=\widetilde t}
 \text{ 中的 }B}_{\text{反事实结构因果模型}},
 \tag{6.27}
\end{equation}
其中 $\mathbf n_u$ 表征单位 $u$［Pearl, 2009，式 (3.51)］。在反事实结构因果模型中，所有噪声项均为确定值，因此导出的 $B$ 的分布也退化，$B$ 也成为确定值，符合要求。表 6.2 的例子利用 Bernoulli 分布 $N_T\sim\operatorname{Ber}(0.6)$ 和 $N_B\sim\operatorname{Ber}(0.01)$，独立同分布地抽取了 200 个单位。此时 SUTVA 成立：独立同分布假设意味着各单位互不干扰；模块性，即干预 $T$ 只改变 $T$ 的结构赋值，则意味着接受治疗的方式不影响结果。

下面讨论一个结果，说明式 (6.27) 中的两种表示在哪种意义下等价。这里主要沿用 Pearl［2009，第 7.3.1 节］和 Halpern［2000］的表述。主要论证步骤如下。
\begin{enumerate}[label=\arabic*.]
 \item 定义性质或公理 (C0)—(C5) 与 (MP)［Halpern, 2000，第 3 节］。例如，性质 (C4) 规定
 \[
 T_u(t=\widetilde t,w=\widetilde w)=t;
 \]
 它要求将单位 $u$ 的变量 $T$ 设为 $t$ 这一操作是“有效的”。
 \item 这些公理在两种表示中都成立，即具有“可靠性”（soundness）。
 \item 可以证明，这些性质对于反事实结构因果模型是完备的。任何反事实陈述，都可由其中一条公理推出。
 \item 因此，对于反事实结构因果模型成立的任意定理，在潜在结局框架中也成立，反之亦然。\footnote{严格地说，“反之亦然”要求潜在结局框架没有在上述公理之外再作额外假设。}
 此外，由第 3 步可知，任何满足这三条公理的数据集，例如表 6.2，都可以用某个反事实结构因果模型建模。\footnote{如果没有任何结构因果模型能够生成该数据集，就意味着结构因果模型中的反事实还满足一个不能由这三条公理推出的额外性质，即“这个数据集不可能被生成”。}
\end{enumerate}
\begin{quote}\small
\textbf{译注：}此处原文有几处简写或符号不一致：第 1 步公式左侧设定 $t=\widetilde t$，右侧应相应读作 $\widetilde t$；第 3 步的完备性，是指所有在相应模型类中\emph{有效}的反事实陈述都可由公理体系推出，并非任意陈述都可由单条公理推出；第 4 步及脚注所称“三条公理”，应理解为前面列出的相关公理体系。原式和原文计数均予保留。
\end{quote}

不过，两种框架所用的语言不同。即使每个定理都在两种框架中成立，有些定理在其中一种框架中可能更“容易”证明。同样，定理中的任何假设，都会限制背后的数据生成过程；根据具体应用，某种表述可能更便于评估这些限制。平均因果效应为零、个体因果效应却非零的情形，似乎更适合用潜在结局处理；另一方面，结构因果模型的图表示，可能更便于利用随机变量之间因果关系的假设。

Richardson and Robins［2013］提出使用\textbf{单世界干预图（single world intervention graph）}。这类图允许将变量设为特定值，从而为反事实变量构造图上的对应表示。修改后的图允许我们读出同时涉及事实变量与反事实变量的条件独立陈述。因此，可以将其看作一种有用工具，将图上的假设转化为潜在结局研究者常用的反事实陈述。
\FloatBarrier
